% Auto-generated by src/make_tables.py -- do not edit by hand.
\begin{table}[t]
\centering
\caption{The reconstruction order in the currency it trades in. Each instrument is run through the executed closed loop with the paper's real-market calibration ($\eta$ = half the Roll effective spread, $\sigma$ from minute bars), with the exact value function placed on the grid so the solve contributes nothing and only the reconstruction differs. \emph{Excess over exact} is the cost of an arm measured against the exact control, in basis points of notional; \emph{multilinear vs cubic} is the pairwise difference, which is the well-posed statement because both arms share the snapshots and the time interpolation. The cubic arm sits at the numerical floor ($\le 0.02$\,bp), so raising the reconstruction order is worth the difference column --- and it is worth most on the MOST liquid instrument, because a tighter spread buys a larger optimal control and therefore a larger reconstruction error. Full horizon: truncating the integral at $T-\delta$ removes the terminal penalty that pays for not liquidating, which makes the slower policy look cheap (the pairwise figure falls to $-0.21$\,bp at $\delta = 0.1$ on SPY); the sensitivity is in the reproduction.}
\label{tab:shortfall}
\scriptsize
\setlength{\tabcolsep}{4pt}
\begin{tabular}{lrrr}
\toprule
instrument & $\eta$ (bp) & excess over exact: ml & ml vs cubic \\
\midrule
SPY & 1.28 & 8.42 & 8.40 \\
AMZN & 1.46 & 7.45 & 7.43 \\
MSFT & 1.58 & 6.91 & 6.89 \\
QQQ & 1.60 & 6.82 & 6.81 \\
GOOGL & 1.97 & 5.60 & 5.59 \\
OXY & 2.24 & 4.97 & 4.96 \\
NVDA & 2.68 & 4.21 & 4.20 \\
AAPL & 2.73 & 4.16 & 4.14 \\
TSLA & 4.41 & 2.69 & 2.67 \\
\midrule
\multicolumn{4}{l}{\emph{mean} 5.68\,bp of notional} \\
\bottomrule
\end{tabular}
\end{table}
